% Options for packages loaded elsewhere
\PassOptionsToPackage{unicode}{hyperref}
\PassOptionsToPackage{hyphens}{url}
\documentclass[
]{article}
\usepackage{xcolor}
\usepackage{amsmath,amssymb}
\setcounter{secnumdepth}{-\maxdimen} % remove section numbering
\usepackage{iftex}
\ifPDFTeX
  \usepackage[T1]{fontenc}
  \usepackage[utf8]{inputenc}
  \usepackage{textcomp} % provide euro and other symbols
\else % if luatex or xetex
  \usepackage{unicode-math} % this also loads fontspec
  \defaultfontfeatures{Scale=MatchLowercase}
  \defaultfontfeatures[\rmfamily]{Ligatures=TeX,Scale=1}
\fi
\usepackage{lmodern}
\ifPDFTeX\else
  % xetex/luatex font selection
\fi
% Use upquote if available, for straight quotes in verbatim environments
\IfFileExists{upquote.sty}{\usepackage{upquote}}{}
\IfFileExists{microtype.sty}{% use microtype if available
  \usepackage[]{microtype}
  \UseMicrotypeSet[protrusion]{basicmath} % disable protrusion for tt fonts
}{}
\makeatletter
\@ifundefined{KOMAClassName}{% if non-KOMA class
  \IfFileExists{parskip.sty}{%
    \usepackage{parskip}
  }{% else
    \setlength{\parindent}{0pt}
    \setlength{\parskip}{6pt plus 2pt minus 1pt}}
}{% if KOMA class
  \KOMAoptions{parskip=half}}
\makeatother
\usepackage{longtable,booktabs,array}
\usepackage{caption}
\captionsetup[table]{skip=6pt}
\newcounter{none} % for unnumbered tables
\usepackage{calc} % for calculating minipage widths
% Correct order of tables after \paragraph or \subparagraph
\usepackage{etoolbox}
\makeatletter
\patchcmd\longtable{\par}{\if@noskipsec\mbox{}\fi\par}{}{}
\makeatother
% Allow footnotes in longtable head/foot
\IfFileExists{footnotehyper.sty}{\usepackage{footnotehyper}}{\usepackage{footnote}}
\makesavenoteenv{longtable}
\setlength{\emergencystretch}{3em} % prevent overfull lines
\providecommand{\tightlist}{%
  \setlength{\itemsep}{0pt}\setlength{\parskip}{0pt}}
\usepackage{bookmark}
\IfFileExists{xurl.sty}{\usepackage{xurl}}{} % add URL line breaks if available
\urlstyle{same}
% fallback for those not using the hyperref driver hyperxmp:
\makeatletter
\@ifundefined{xmpquote}{\newcommand{\xmpquote}[1]{#1}}{}
\makeatother
\hypersetup{
  hidelinks,
  pdfcreator={LaTeX via pandoc}}

\author{}
\date{}

\begin{document}

\section{Paper IV v1.22-r2 --- Informe final consolidado de auditoría
externa adversarial
(E0--E8)}\label{paper-iv-v122-r2--informe-final-consolidado-de-auditoruxeda-externa-adversarial-e0e8}

Ejecución: \texttt{run\_v1.22\_r2}. Estado: \textbf{completada} el
2026-10-01: inicio 11:00:18 UTC (E0 inicial 11:01:08 UTC); cierre tras
el E0 final de la §11. Fue solicitada por el propietario en el chat
(«haz la rerevision») con el mandato
\texttt{EXTERNAL\_ADVERSARIAL\_REVALIDATION\_v1.22\_r2.md}.

\textbf{Veredicto global actual: PASS\_WITH\_OBSERVATIONS.} No queda
ningún hallazgo de severidad MINOR o superior abierto. Siguen vigentes
nueve observaciones aceptadas (ocho históricas y la nueva NEW-04, §8),
una mitigación (X-28) y los límites de independencia de la §9. No es una
revisión por pares humana ni autoriza publicar.

\begin{center}\rule{0.5\linewidth}{0.5pt}\end{center}

\subsection{1. Identidad exacta del paquete
revisado}\label{1-identidad-exacta-del-paquete-revisado}

{\def\LTcaptype{none} % do not increment counter
\begin{longtable}[]{@{}lll@{}}
\toprule\noalign{}
Objeto & Ruta & SHA-256 \\
\midrule\noalign{}
\endhead
\bottomrule\noalign{}
\endlastfoot
Objetivo r2 &
\texttt{02\_validation/02\_IA\_ADVERSARIAL\_AUDITS/AUDIT\_TARGET\_v1.22\_r2.json}
&
\texttt{f931776370dcc101c637cb978273d8757b298ff79e99d50fff815104d4e00583}
(coincide con el sidecar) \\
Mandato r2 &
\texttt{…/EXTERNAL\_ADVERSARIAL\_REVALIDATION\_v1.22\_r2.md} &
\texttt{d085621334d1f5ae14f9e637d237817d7d0fe6673709cfd5fcbad6f4c984ba3e} \\
Manuscrito & versión \textbf{v1.22}, paquete \textbf{r2}:
\texttt{01\_manuscript/v1.22\_editorial\_candidate\_r2/} & --- \\
Manifiesto r2 & \texttt{MANUSCRIPT\_REVIEW\_MANIFEST.json} (127
archivos) &
\texttt{0fcaf0589c389f70fdb85959f4c69c45f97a131a64a7234a78e40328f7509c0a} \\
ZIP r2 & \texttt{PAPER\_IV\_v1.22\_REVIEW\_PACKAGE\_r2.zip} (129
miembros, CRC correcto) &
\texttt{d4f3b33d647744a54c93098e0944b18d01a3377c4bbc9a3ecc5f8789e9973635} \\
ES MD / TeX / PDF (73 p.) & \texttt{PAPER\_IV\_preprint\_v1.22\_es.*} &
\texttt{ec227eb5…e0e97e} / \texttt{356f03c1…222d02c} /
\texttt{8ebfccd7…6eb0df} \\
EN MD / TeX / PDF (72 p.) & \texttt{PAPER\_IV\_preprint\_v1.22\_en.*}
--- \textbf{byte-idénticos a r1} & \texttt{99a1159d…65c94f} /
\texttt{34a3081b…9c1100} / \texttt{5bbb730e…cd79e0} \\
Corte Lean (freeze) &
\texttt{05\_formalization/lean\_piv-v12-fb459343d234/}, 615 entradas &
manifiesto
\texttt{fb459343d234f968d7d32eff1491ea8a09aa2e135b313a80623012e7449042f5} \\
ZIP de fuentes Lean &
\texttt{05\_formalization/LEAN\_SOURCE\_piv-v12-fb459343d234.zip} (615
miembros = manifiesto) &
\texttt{cb2741454380736ffe01b59933057b5079fa9f865d5bbdd3f67534bf808a62e6} \\
Anexo & \texttt{05\_formalization/LEAN\_BOUNDED\_GAP\_ANNEX\_v1.0.zip}
(40 miembros) &
\texttt{2847a422865e06880d457ea806aae5b9f749d22359eff325349c09c01cb4837d} \\
Corte r1 de comparación &
\texttt{01\_manuscript/v1.22\_editorial\_candidate/} (sin cambios) &
manifiesto \texttt{d9e559a3…2984029a}; ZIP \texttt{2b87b746…f9e2b098} \\
Paquetes históricos &
\texttt{run\_v1.2\_r1/40\_PACKAGE/EXTERNAL\_AUDIT\_run\_v1.2\_r1.zip} &
\texttt{886ed7f033bca42a530944d48f91ef45e79f6be4cb0d74080c5d8d6efc4f5e2b} \\
&
\texttt{run\_v1.21\_r1/40\_PACKAGE/EXTERNAL\_REVALIDATION\_run\_v1.21\_r1.zip}
&
\texttt{11659a97b0a4f2bffadb6968159cc5b66562b4470053939a45c1e8a42db7a638} \\
&
\texttt{run\_v1.22\_r1/40\_PACKAGE/EXTERNAL\_REVALIDATION\_run\_v1.22\_r1.zip}
&
\texttt{9c74d673206ab17a5b4586209ececca30536f1768fcc91ffdd39db551193a1fa} \\
\end{longtable}
}

Los hashes completos del manuscrito están en
\texttt{20\_EVIDENCE/E0/E0\_INITIAL.json}. E0 inicial y E0 final dan
resultados idénticos (§11). Ningún insumo ni evidencia histórica cambió.

\subsection{2. Tabla E0--E8: veredicto
actual}\label{2-tabla-e0e8-veredicto-actual}

{\def\LTcaptype{none} % do not increment counter
\begin{longtable}[]{@{}llllll@{}}
\toprule\noalign{}
Puerta & Veredicto actual & Alcance & Base del veredicto & Procedencia
(ejecución, fecha) & Evidencia (ruta / hash) \\
\midrule\noalign{}
\endhead
\bottomrule\noalign{}
\endlastfoot
E0 identidad & \textbf{PASS} & objetivo, manifiestos, ZIP, inglés, Lean,
anexo, paquetes históricos & \textbf{ejecutado en r2} (inicial y final)
& run\_v1.22\_r2, 2026-10-01 &
\texttt{20\_EVIDENCE/E0/E0\_INITIAL.json}, \texttt{E0\_FINAL.json} \\
E1 enunciados & \textbf{PASS} & 21 afirmaciones del mapa; constantes y
etiquetas & \textbf{PASS heredado y revalidado por identidad}, más un
control de impacto nuevo & run\_v1.2\_r1 (2026-09-30), run\_v1.21\_r1,
run\_v1.22\_r1 (2026-10-01) &
\texttt{run\_v1.2\_r1/20\_EVIDENCE/E1/E1\_RECORD.md};
\texttt{20\_EVIDENCE/IMPACT/IMPACT\_CHECK.json} \\
E2 matemática & \textbf{PASS} (observación R-01) & 56 componentes; 7
filas de A.2 & \textbf{PASS heredado y revalidado por identidad} &
A2-1/2/7: run\_v1.2\_r1; A2-3/5/6: run\_v1.21\_r1; A2-4: run\_v1.22\_r1
& \texttt{run\_v1.2\_r1/…/E2\_RECORD.md};
\texttt{run\_v1.21\_r1/…/E2\_RECORD.md};
\texttt{run\_v1.22\_r1/…/E2\_RECORD.md} \\
E3 correspondencia & \textbf{PASS} & cabeceras Lean, identificadores,
X-27 & \textbf{PASS heredado y revalidado por identidad}, más la
resolución de \texttt{Model} & run\_v1.2\_r1 → run\_v1.22\_r1 &
\texttt{20\_EVIDENCE/E3/E3\_RECORD.md} \\
E4 build & \textbf{PASS --- reused\_verified\_external\_build} & 607
módulos, 19 targets, 224 exportaciones, 50 módulos del anexo &
\textbf{PASS heredado y revalidado por identidad}; \textbf{sin build ni
replay del kernel} & build de run\_v1.2\_r1, 2026-09-30 15:09--21:33 UTC
& consola \texttt{ef5b01c7…}, records \texttt{0f2d8495…}, anexo
\texttt{57ce1895…}; \texttt{20\_EVIDENCE/E4/E4\_REUSE\_CHECK.json} \\
E5 falsación & \textbf{PASS} & T01--T24; 21 + 14 comprobaciones exactas
& \textbf{PASS heredado y revalidado por identidad}, con repetición de
los scripts aritméticos & run\_v1.2\_r1, run\_v1.21\_r1, run\_v1.22\_r1
& \texttt{20\_EVIDENCE/E5/*\_RERUN.json} \\
E6 bilingüe / PDF & \textbf{PASS\_WITH\_OBSERVATIONS} & delta ES r1→r2,
73 páginas, PDF frente a TeX & \textbf{ejecutado en r2} & run\_v1.22\_r2
& \texttt{20\_EVIDENCE/E6/E6\_RECORD.md}, \texttt{E6\_R2\_DELTA.json},
\texttt{PAGE\_INSPECTION\_LOG.csv} \\
E7 literatura & \textbf{PASS} & bibliografía, atribución & \textbf{PASS
heredado y revalidado por identidad} (referencias idénticas) &
run\_v1.2\_r1 → run\_v1.22\_r1 &
\texttt{20\_EVIDENCE/E7/E7\_RECORD.md} \\
E8 evidencia del autor & \textbf{PASS\_WITH\_OBSERVATIONS} & detector,
alcance 1266/844, paquete, NEW-03 & detector \textbf{repetido en r2} y
atribución heredada & run\_v1.21\_r1 (barrido 1266), run\_v1.22\_r2 &
\texttt{20\_EVIDENCE/E8/E8\_RECORD.md},
\texttt{E8\_DETECTOR\_RERUN.json} \\
\end{longtable}
}

La matriz completa (veredicto por revisión, comprobación nueva,
evidencia heredada y enlace de identidad) está en
\texttt{30\_REPORT/TRACEABILITY\_MATRIX.csv}.

\subsection{3. Historial de revisiones}\label{3-historial-de-revisiones}

{\def\LTcaptype{none} % do not increment counter
\begin{longtable}[]{@{}lllll@{}}
\toprule\noalign{}
Revisión & Ejecución & Veredicto global original & Puertas no PASS & Qué
abrió o cerró \\
\midrule\noalign{}
\endhead
\bottomrule\noalign{}
\endlastfoot
v1.2 & run\_v1.2\_r1 (2026-09-30) & \textbf{INCONCLUSIVE} & E2
INCONCLUSIVE (A2-3, A2-4, A2-5 y A2-6: REQUIRES\_EXPANSION); E3, E6 y E8
PASS\_WITH\_FINDINGS; E7 PASS\_WITH\_MINOR\_FINDINGS & Hallazgos
X-01\ldots X-29: 5 MODERATE, 13 MINOR, 11 OBS/INFO. Build externo
607/607, 19/19, 50/50, oleans idénticos byte a byte. X-27 (procedencia
de E32) \\
v1.21 & run\_v1.21\_r1 (2026-09-30/10-01) & \textbf{INCONCLUSIVE} & E2
INCONCLUSIVE (A2-4 seguía abierto); E1 PASS\_WITH\_MINOR; E6
PASS\_WITH\_FINDINGS; E8 PASS\_WITH\_OBSERVATIONS & Cerró A2-3, A2-5 y
A2-6 (X-01, X-02, X-03) y otros 15. Abrió NEW-01, NEW-02, R-01, R-02,
E8-01\ldots03. X-05, X-12 y X-15 quedaron parciales \\
v1.22-r1 & run\_v1.22\_r1 (2026-10-01) & \textbf{PASS} («con hallazgos
MINOR no bloqueantes») & E6 PASS\_WITH\_FINDINGS; E8
PASS\_WITH\_OBSERVATIONS & Cerró A2-4/X-05 (C.3), X-12, NEW-02, R-02,
E8-01, E8-03 y la ortografía de X-10. Dejó parciales X-15 y NEW-01 (ES).
Abrió NEW-03 \\
v1.22-r2 & \textbf{run\_v1.22\_r2} (2026-10-01) &
\textbf{PASS\_WITH\_OBSERVATIONS} & E6 y E8 PASS\_WITH\_OBSERVATIONS &
Cierra X-15, NEW-01 y NEW-03. Registra la observación NEW-04 y la
corrección del auditor C3-01 \\
\end{longtable}
}

Los INCONCLUSIVE de v1.2 y v1.21 \textbf{se mantienen como veredictos de
esas revisiones}. El PASS actual se refiere al paquete v1.22-r2 y no las
convierte retroactivamente.

\subsection{4. Cobertura matemática
consolidada}\label{4-cobertura-matemuxe1tica-consolidada}

\subsubsection{4.1 Afirmaciones
públicas}\label{41-afirmaciones-puxfablicas}

Mapa de afirmaciones de run\_v1.2\_r1, E1 (21 filas). Su texto es
idéntico en r2.

{\def\LTcaptype{none} % do not increment counter
\begin{longtable}[]{@{}lllll@{}}
\toprule\noalign{}
Afirmación & Declaración Lean & Semántica (E1) & Demostración escrita
(E2) & Formal (E3/E4) \\
\midrule\noalign{}
\endhead
\bottomrule\noalign{}
\endlastfoot
Teorema A (todos los órdenes, aditiva) &
\texttt{Erdos81AllOrders.erdos81\_all\_orders\_additive}; explícita
\texttt{ExplicitThreshold.erdos81\_all\_orders\_bounded\_additive} &
MATCH & ensamblaje rederivado; b = N² ≤ T(h+8) vía C.3 (A2-4, cerrado en
r1) & compilado, axiomas estándar \\
Teorema B (1.1) cota y (1.2) máximo &
\texttt{Erdos81Unconditional.erdos81\_cliquePartition};
\texttt{PaperTheorems.erdos81\_max\_eq} & MATCH, MATCH & rederivado
(ramas L/C, Corolario 3.5, Teorema 5.0, Lema 5.3) & idem \\
Teorema C (defecto fijo) &
\texttt{DefectExplicitPublication.rooted\_defect\_eventual/\_maximum} &
MATCH (rsd equivalente a {[}15{]}) & E.1--E.4; A2-5 y A2-6 cerrados en
v1.21 & idem; cono sin Alon--Shapira \\
Teorema C′ (6.12), (6.13); Cor. 6.3 (6.17); clasificación (6.18) &
\texttt{SublinearResearch.fixed\_defect\_joint\_stability\_real},
\texttt{fixed\_defect\_exact\_extremal\_edit\_real},
\texttt{E32.cp\_classification\_of\_theoremCPrimeC} & MATCH (C′a--d) &
rederivado; la misma raíz antes de Q y τ & idem; \textbf{X-27}: la
clasificación cp usa la cadena histórica {[}23{]}, declarado \\
Estabilidad cordal: Teorema 6.1, Cor. 6.1a, 6.2 &
\texttt{IntegralStability…}, \texttt{chordal\_joint\_stability\_real},
\texttt{ExtremalClassification.chordal\_extremal\_classification} &
MATCH & rederivado (componente 26) & idem \\
Prop. 6.3a (raíz cuadrada) & \texttt{OptimalTemplateObstruction.*} &
MATCH; X-17 (Lean existencial) & rederivado & idem \\
Prop. 6.4, Teorema 6.5, Cor. 6.6, (C.4) &
\texttt{certified\_*\_all\_orders}, \texttt{uniform\_*},
\texttt{sequence\_arbitrary\_orders\_*}, \texttt{E35.theoremC\_tower*},
\texttt{E35.NfarE\_le\_tower\_poly} & MATCH & estático y rederivado; la
torre de (C.4) se compara en C.3 & idem \\
Teorema 3.1, Cor. 3.5 & \texttt{RC01Final.rc01\_uniformRoundingTarget},
\texttt{FarRegimeAllGraphs…} & MATCH & C.1--C.3; A2-3 cerrado en v1.21 &
idem \\
Prop. D.2, D.3, D.1 &
\texttt{DefectSharpPublication.order\_three\_insufficient},
\texttt{SplitMixedGap.mixed\_gap\_zero},
\texttt{ThreeRegime.CompleteStateAllOrders} & D.2 MATCH & componentes
52--53 rederivados; A2-7 & idem \\
F.3a/b & \texttt{exists\_clique\_additive\_defect},
\texttt{recover\_clique\_from\_edit} & MATCH & rederivado & idem \\
Anexo G.1 & \texttt{BoundedCliqueGap.chordal\_gap\_linear\_cliqueFree} &
MATCH & --- & anexo 50/50, 633 declaraciones, axiomas estándar \\
G.2 (Prop. G.1) &
\texttt{FarExploration.CleanupRigidVerdict.threshold\_gt\_exp} &
\textbf{no re-extraída} (solo texto estático) & --- & compilado dentro
del corte \\
\end{longtable}
}

\textbf{Diferencias declaradas.} La Tabla 7 del manuscrito incluye
entradas que \textbf{no tienen fila semántica propia en E1}: Teorema 5.0
(ventana reforzada e interfaz anterior), Cor. 5.4, Dicotomía (1.5),
optimalidad de los coeficientes cuadrático y lineal, identidad (6.7b) y
Prop. D.3. Están cubiertas solo por el registro de componentes de E2
(componentes 14, 22, 24, 26 y 53, rederivados) y por la correspondencia
estática de E3 más la compilación de E4. No les atribuyo una cobertura
semántica que no existe. G.2 sigue sin re-extraer.

\subsubsection{4.2 Las siete derivaciones de
A.2}\label{42-las-siete-derivaciones-de-a2}

{\def\LTcaptype{none} % do not increment counter
\begin{longtable}[]{@{}llllll@{}}
\toprule\noalign{}
ID auditor & Fila literal de A.2 en el manuscrito (ES p. 41 / EN p. 40)
& v1.2 (run\_v1.2\_r1) & v1.21 & v1.22-r1 & Estado actual \\
\midrule\noalign{}
\endhead
\bottomrule\noalign{}
\endlastfoot
A2-1 & «§5.1, (5.4) y (5.7)» & ACCEPTABLE\_SUMMARY & idem (X-23: frase
113/64 añadida) & heredado & \textbf{ACCEPTABLE\_SUMMARY} \\
A2-2 & «C.2, (3.8) y (3.10)» & ACCEPTABLE\_SUMMARY & idem & heredado &
\textbf{ACCEPTABLE\_SUMMARY} \\
A2-3 & «Lema 3.2 y C.1--C.2» & REQUIRES\_EXPANSION (la carga por arista
no se seguía del texto) & \textbf{ACCEPTABLE\_SUMMARY} (limpieza
bilateral, carga ≤ 1, retención; X-01) & heredado; R-02 aclarado &
\textbf{ACCEPTABLE\_SUMMARY} (R-01 obs.) \\
A2-4 & «C.3, Tabla C.1» & REQUIRES\_EXPANSION & REQUIRES\_EXPANSION
(faltaba la recurrencia de regularidad y la comparación con la torre) &
\textbf{ACCEPTABLE\_SUMMARY}: 6 comprobaciones contra Mathlib
\texttt{initialBound}/\texttt{stepBound}/\texttt{bound} y E18/E19 &
\textbf{ACCEPTABLE\_SUMMARY} \\
A2-5 & «E.1» (normalización) & REQUIRES\_EXPANSION &
\textbf{ACCEPTABLE\_SUMMARY} (X-02) & heredado &
\textbf{ACCEPTABLE\_SUMMARY} \\
A2-6 & «E.1, (E.2b) y (E.2d)» & REQUIRES\_EXPANSION &
\textbf{ACCEPTABLE\_SUMMARY} (X-03) & heredado &
\textbf{ACCEPTABLE\_SUMMARY} \\
A2-7 & «D.1» & ACCEPTABLE\_SUMMARY & idem & heredado &
\textbf{ACCEPTABLE\_SUMMARY} \\
\end{longtable}
}

Diferencias entre las filas literales y los ID del auditor:

\begin{itemize}
\tightlist
\item
  A2-1: el auditor citaba también (5.3).
\item
  A2-3: el auditor la describía como «§3.2/C.2, cotas de fibra, (3.8),
  factibilidad por arista», mientras que el manuscrito agrupa (3.8) en
  la fila C.2.
\item
  A2-5: el auditor incluía además las cuentas de E.4.
\end{itemize}

El contenido cubierto es el mismo, pero esas son las diferencias de
redacción.

\subsection{5. E4: build externo y por qué sigue
valiendo}\label{5-e4-build-externo-y-por-quuxe9-sigue-valiendo}

\textbf{Qué se hizo} (run\_v1.2\_r1, 2026-09-30):

\begin{itemize}
\tightlist
\item
  Runner propio del auditor, sin Lake, con Lean 4.28.0. Fuentes
  extraídas del ZIP (615/615) y objetos nuevos, sin usar oleans del
  autor. Mathlib compartido en solo lectura; caché de 113 821 archivos
  sin cambios.
\item
  El primer intento falló por un error del auditor (C-04); el segundo
  corrió de 15:09 a 21:33 UTC.
\end{itemize}

\textbf{Registros originales:}

\begin{itemize}
\tightlist
\item
  \textbf{Consola:} \texttt{modules\_planned\ 607,\ pass\ 607}, sin
  fallos, 607 líneas \texttt{PASS\ n\ 607}, \texttt{EXIT\ 0}, y los 19
  targets de FREEZE\_SCOPE en PASS.
\item
  \textbf{records.jsonl:} 608 filas (607 + AuditorChecks), todas PASS,
  con código de salida 0 y hashes de fuente iguales al manifiesto.
\item
  \textbf{Logs:} 608 logs de módulo sin «sorry».
  \texttt{ReleaseExportCheck} tiene 224 \texttt{\#check} y 0 errores.
\item
  \textbf{Anexo:} 50/50 PASS; 633 declaraciones auditadas con axiomas
  estándar.
\item
  \textbf{Identidad:} los 607 oleans son idénticos byte a byte a los del
  autor.
\end{itemize}

\textbf{Huella de axiomas.} No son teoremas distintos:

\begin{itemize}
\tightlist
\item
  461 \emph{registros} de axiomas según la regla del autor, de los
  cuales 314 son salidas de \texttt{\#print\ axioms} y el resto líneas
  de auditoría repetidas.
\item
  18 declaraciones principales con \texttt{\#print\ axioms} = \{propext,
  Classical.choice, Quot.sound\}.
\item
  Conos propios de 50 declaraciones públicas: 0 axiomas no estándar, 0
  \texttt{sorryAx}, 0 constantes \texttt{Erdos81}.
\end{itemize}

\texttt{E4\_main/SUMMARY.json} fue sobrescrito por la invocación de
AuditorChecks: ahora dice «555 planned» y solo lista AuditorChecks.
\textbf{No se usa} como evidencia.

\textbf{Por qué sigue valiendo para r2.}

\begin{itemize}
\tightlist
\item
  r2 no toca Lean.
\item
  E0 confirma, al inicio y al final, las 615/615 entradas, el ZIP de
  fuentes igual al manifiesto, el anexo y los registros de run\_v1.2\_r1
  sin cambios.
\item
  El script documental (idéntico al de r1) reproduce exactamente el
  mismo resultado.
\end{itemize}

\textbf{No se ejecutó en r2 ningún build nuevo ni un replay
independiente del kernel.}

\subsection{6. E5 y E8: pruebas, controles y
recuentos}\label{6-e5-y-e8-pruebas-controles-y-recuentos}

\textbf{E5.}

\begin{itemize}
\tightlist
\item
  Inventario de comprobaciones:

  \begin{itemize}
  \tightlist
  \item
    T01--T24 de run\_v1.2\_r1: 24 pruebas con controles negativos
    predeclarados. La verificación del certificado Certo confirma una
    partición exacta de K3∨I3 en 6 piezas, solo como factibilidad.
  \item
    21 comprobaciones exactas de v1.21.
  \item
    14 de C.3 (r1), más el control negativo corregido de C3m y 6
    controles negativos complementarios.
  \end{itemize}
\item
  En total, 59 comprobaciones distintas y 7 controles añadidos. Las
  repeticiones de r2 dan resultados idénticos y \textbf{no se suman}
  como pruebas nuevas.
\item
  Las comprobaciones sobre rangos finitos son regresión, no prueba
  universal; la prueba universal está en E2 y en E19.
\end{itemize}

\textbf{E8.}

\begin{itemize}
\tightlist
\item
  El barrido de 1266 logs = 608 + 50 (auditor) + 422 + 186 (segmentos
  del autor) es del auditor de run\_v1.21\_r1. El barrido antiguo del
  autor cubría 844 = 1266 − 422 y omitía el segmento 1 (E8-01).
\item
  r2 repitió el mismo barrido con resultado idéntico (0 alertas). Cuenta
  una sola vez.
\item
  \textbf{X-28 sigue MITIGADO:} los runners congelados mantienen un
  patrón insuficiente, su hash no cambió y el parche no se aplicó. La
  mitigación es el detector externo (10/10 casos) más los conos de
  axiomas.
\item
  \textbf{R-01 sigue como limitación:} las cotas de segundo momento de
  C.2 se citan.
\end{itemize}

\subsection{7. Comprobaciones nuevas de
r2}\label{7-comprobaciones-nuevas-de-r2}

\begin{itemize}
\tightlist
\item
  \textbf{X-15.} La prosa española de r1 tenía 11 «matching(s)» y 1
  «packing»; en r2 no queda ninguno.

  \begin{itemize}
  \tightlist
  \item
    Las 4 apariciones que permanecen son títulos citados ({[}6{]},
    {[}7{]}, {[}9{]}, {[}12{]}) y no se tradujeron, como corresponde. No
    hay ninguna en identificadores.
  \item
    \textbf{Discrepancia de recuento resuelta:} el «×12 matching(s)» de
    los informes anteriores era un error del auditor. Contaba la
    referencia {[}9{]}, y el escaneo no detectaba la forma plural. El
    recuento correcto es 11 + 1 (CORRECTIONS C3-01).
  \item
    \textbf{CERRADO.}
  \end{itemize}
\item
  \textbf{NEW-01.} \texttt{Model} en monoespaciado en §7.2: MD l.1203,
  TeX l.1084 y PDF p. 35 (140 dpi). \textbf{CERRADO.}
\item
  \textbf{Diff completo.} Mi diff coincide byte a byte con
  \texttt{CHANGES\_es.diff}. Hay 13 cambios de palabra en 11 líneas,
  idénticos en MD y TeX.

  \begin{itemize}
  \tightlist
  \item
    Display (237), en línea (1989), bloques de código (5), etiquetas
    (146) y referencias: idénticos.
  \item
    Ninguna línea cambiada es un enunciado; una está dentro de una
    demostración (Lema 5.3).
  \end{itemize}
\item
  \textbf{El PDF corresponde al TeX.} El texto del PDF r2 coincide con
  el de r1 tras las sustituciones del inventario, salvo 3 diferencias
  tipográficas verificadas visualmente:

  \begin{itemize}
  \tightlist
  \item
    espaciado de una línea matemática en p. 55;
  \item
    la Tabla 9, que ahora cabe en una sola página.
  \item
    El log está limpio (0 Overfull, solo 5 Underfull) y el preámbulo es
    idéntico.
  \end{itemize}
\item
  \textbf{Páginas.}

  \begin{itemize}
  \tightlist
  \item
    Cambian 18 páginas en raster (16, 18, 20, 21, 35, 55, 56, 61--71);
    coincide con la lista del editor.
  \item
    Revisé las 73 páginas. Las páginas 55, 56 y 61--71 se compararon
    además lado a lado r1\textbar r2, revisando los desplazamientos; no
    hay texto perdido ni duplicado, y las referencias siguen empezando
    en p. 72.
  \item
    Las comparaciones r1\textbar r2 de las páginas 16, 18, 20, 21 y 35
    no llegaron a mostrarse (C3-02). Esas páginas se verificaron en las
    hojas r2 y, la 35, también a 140 dpi.
  \end{itemize}
\item
  \textbf{NEW-03.} La rectificación explícita del editor y el artefacto
  corregido cierran el hallazgo. La respuesta r1 y el hallazgo histórico
  se conservan. \textbf{CERRADO.}
\item
  \textbf{Impacto en E1, E2, E3 y E7.} Ninguna conclusión previa cambia
  (§2).
\end{itemize}

\subsection{8. Hallazgos: lista completa y estado
actual}\label{8-hallazgos-lista-completa-y-estado-actual}

Hay 49 filas en \texttt{FINDINGS.csv}: todas las históricas conservadas
más 4 nuevas.

\begin{itemize}
\tightlist
\item
  \textbf{Cerrados (27):}

  \begin{itemize}
  \tightlist
  \item
    X-01, X-02, X-03 y X-04 (MODERATE), cerrados en v1.21.
  \item
    X-05, cerrado en v1.22-r1.
  \item
    X-06\ldots X-11 y X-13, X-14, X-16 (MINOR), cerrados en v1.21; X-12,
    en v1.22-r1.
  \item
    X-15, en v1.22-r2.
  \item
    X-20, X-21, X-23 y X-25 (OBS).
  \item
    X-27 (MODERATE), cerrado por declaración en v1.21: el código Lean no
    cambió y la declaración se conserva.
  \item
    NEW-01 (v1.22-r2), NEW-02 (v1.22-r1), NEW-03 (v1.22-r2), R-02
    (v1.22-r1), E8-01 y E8-03 (v1.22-r1).
  \item
    Responsables: el editor o autor corrigió el texto; el auditor
    verificó en la revisión indicada.
  \end{itemize}
\item
  \textbf{Mitigado (1):} X-28, el detector de los runners congelados.
\item
  \textbf{Observaciones aceptadas (8 vigentes + NEW-04):}

  \begin{itemize}
  \tightlist
  \item
    X-17 (Prop. 6.3a, existencial en Lean) y X-18 (τ real).
  \item
    X-19 (residuo del lakefile) y X-22 (desacuerdo interno/externo sobre
    A.2, conservado).
  \item
    X-24 (constantes opcionales no adoptadas) y X-29 (reproducción en
    dos tramos del autor).
  \item
    R-01 (segundo momento citado) y E8-02 (alcance del parche).
  \item
    \textbf{NEW-04 (nueva).} El texto histórico de estado del manuscrito
    (ES/EN p. 1, §7 «Estado de auditoría», nota de la Tabla 7 y alcance
    de A.2) dice que las ampliaciones de v1.22 «esperan revalidación».
    Responsable: el editor, antes de publicar.
  \end{itemize}
\item
  \textbf{Información:} X-26, sobre independencia.
\item
  \textbf{Correcciones conservadas (11):}

  \begin{itemize}
  \tightlist
  \item
    del proceso v1.2: AUD-P1, AUD-C1 y AUD-C2;
  \item
    del auditor en r1: C2-01 a C2-08;
  \item
    del auditor en r2: C3-01 (recuento), C3-02 (filas de inspección no
    vistas, descartadas) y C3-03 (ruta de salida de un script).
  \end{itemize}
\end{itemize}

\subsection{9. Independencia y
límites}\label{9-independencia-y-luxedmites}

\begin{itemize}
\tightlist
\item
  \textbf{Auditor:} Claude Opus 5.5 (\texttt{claude-opus-5-5}),
  proveedor Anthropic.
\item
  \textbf{Sesión:} esta revalidación es \textbf{continuación de la misma
  sesión de Claude Code} que produjo run\_v1.2\_r1, run\_v1.21\_r1 y
  run\_v1.22\_r1, con compactación de contexto. \textbf{No es una sesión
  nueva, no es ciega y no es independiente de familia.} El manuscrito
  declara el uso de Claude (Anthropic): misma familia.
\item
  \textbf{Lecturas:} leí la evidencia previa y los documentos del
  editor, contrastándolos con la mía. Las lecturas están en
  \texttt{00\_CONTROL/INPUT\_ACCESS\_LOG.csv}. La restricción de lectura
  es operativa, no un aislamiento técnico.
\item
  \textbf{Entorno:} misma máquina y misma caché Mathlib compartida que
  en el build histórico. En r2 no hubo ejecución dinámica nueva de Lean:
  sin build, sin replay del kernel y sin leanchecker. Solo scripts
  Python de solo lectura.
\item
  \textbf{Alcance de la revisión escrita:} la comprobación de las
  demostraciones es a nivel de resumen ejecutable frente a enunciados
  formales congelados. \textbf{No sustituye la revisión humana por
  pares.}
\item
  \textbf{Inspección visual:} hojas a 90 dpi, comparaciones a 110 dpi y
  zoom a 140 dpi. Algunas imágenes no llegaron a mostrarse, y esas filas
  se descartaron (C3-02).
\end{itemize}

\subsection{10. Veredicto razonado y pendientes
reales}\label{10-veredicto-razonado-y-pendientes-reales}

\textbf{PASS\_WITH\_OBSERVATIONS} para el paquete v1.22-r2.

\textbf{Qué certifica esta revisión consolidada:}

\begin{itemize}
\tightlist
\item
  la identidad exacta del paquete revisado;
\item
  que las correcciones españolas pedidas están aplicadas y no alteran el
  contenido protegido;
\item
  que el PDF corresponde al TeX;
\item
  que el inglés y el corte Lean son los mismos que se auditaron;
\item
  que los PASS de E1--E5 y E7 siguen enlazados por identidad a su
  evidencia original;
\item
  que no queda ningún hallazgo MINOR o superior abierto.
\end{itemize}

\textbf{Qué no certifica:}

\begin{itemize}
\tightlist
\item
  no es una revisión por pares humana ni una verificación independiente
  de familia o sesión;
\item
  no hay un build nuevo ni un replay del kernel;
\item
  no prueba la novedad universal;
\item
  no autoriza publicar;
\item
  no cubre semánticamente, fila por fila, las entradas de la Tabla 7
  enumeradas en §4.1.
\end{itemize}

\textbf{Pendientes reales:}

\begin{enumerate}
\def\labelenumi{\arabic{enumi}.}
\tightlist
\item
  \textbf{Antes de publicar}, actualizar el texto histórico de estado en
  ambos idiomas (NEW-04). Esto rompe la identidad byte a byte del inglés
  y exige un control de identidad y un E6 nuevos.
\item
  Las observaciones aceptadas siguen vigentes (R-01, X-28, X-17, X-18 y
  demás).
\item
  Una revisión por pares humana y, si se quiere independencia real, una
  auditoría de otra familia de modelos en una sesión separada, con un
  build de sala limpia o un replay del kernel.
\end{enumerate}

\textbf{Texto histórico de estado.}

\begin{itemize}
\tightlist
\item
  Las menciones de los manuscritos («la revalidación de 1.21\ldots»,
  «esperan revalidación», «espera una nueva revisión») describen el
  estado anterior a run\_v1.22\_r1. No son el estado actual del paquete;
  el estado actual es el de este informe.
\item
  Mientras el paquete sea un candidato de revisión, no causan un error
  matemático ni de correspondencia.
\item
  \textbf{Sí son un problema editorial real si se publica tal cual},
  porque el lector creería que C.3 no está revisado.
\end{itemize}

\subsection{11. E0 final y sellado}\label{11-e0-final-y-sellado}

\begin{itemize}
\tightlist
\item
  \textbf{E0 final:} \texttt{20\_EVIDENCE/E0/E0\_FINAL.json}. Su objeto
  \texttt{checks} es idéntico al del inicial (resultado detallado en
  \texttt{20\_EVIDENCE/logs/E0\_final.log}).
\item
  \textbf{Paquete:}
  \texttt{40\_PACKAGE/EXTERNAL\_REVALIDATION\_run\_v1.22\_r2.zip}, con
  este informe, su PDF, SUMMARY, FINDINGS, la matriz, los registros, los
  scripts, las salidas, las imágenes de E6 y
  \texttt{RUN\_MANIFEST.json}. Los ZIP históricos no se duplican: se
  enlazan por ruta y hash (§1).
\item
  \textbf{Generación del PDF:} pandoc (Markdown → HTML y TeX
  independiente) más PyMuPDF Story (HTML → PDF). \textbf{No se compila
  desde TeX}, el mismo método declarado en las ejecuciones anteriores.
  El PDF se revisó visualmente.
\item
  \textbf{Sin efectos externos:} nada se publicó, se hizo push ni se
  envió a terceros. No se modificaron el manuscrito, los originales ni
  las revisiones anteriores.
\end{itemize}

\end{document}
